\documentclass[12pt,a4paper]{report}

\input{preamble_PPS.tex}

\title{Preliminary Problem Statement: Concurrency analysis of LLM generated programs.}
\author{Tore Asbjørn Tonnisen Kjelds \\ tokj@itu.dk}
\date{\today}
\begin{document}

\maketitle

\section{Context}
AI usage in software engineering has become an integral part of writing software, with one survey claiming more than 90\% of developers use AI tools \cite{jetbrains2025developer}. Recent developments in agentic harnesses (Example: Copilot, Codex, Claude code and Cursor) have also been widely adopted with the same surveyed group, about 29\% use Github Copilot for their work.\newline 
The use of concurrency is motivated by the potential performance gains enabled by modern multicore-CPUs. However, concurrency introduces significant complexity. Managing multiple threads that share ressources can create a seemingly stochastic execution environment in which repeated runs of the smae program produce different results. This, combined with the stochastic nature LLM's 

\section{State of the research}
Many current benchmarks and analysis of LLM generated code focus primarily on sequential programs, leaving a blindspot for concurrent programs. The Concur benchmark \cite{huang2026concur} was developed in part to address this gap. It contains a base set of more than 40 concurrency problems drawn from concurrency textbooks, and uses formal model checking to validate concurrency properties of the generated programs. \newline
The study founc that LLM generated programs vary considerably in their concurrency correctness. However model checking can be very time-consuming, which may restrict use in actual development environments. The study also only examined chat-models, and not agentic harneses, which continues to rise in popularity among developers. \newline
Other studies have investigated scalable concurrency analysis \cite{jåtten2026scalablethreadsafetyanalysisjava}, in which the tool \textbf{CodeQL} was used across some of the most popular open-source repos. The findings of that study resulted in mulitple pull-requests for concurrency fixes. 

\section{Proposed solution}
Given the importance of concurrency in modern software, along with the prevalence of LLM-generated code. It is important to investigate capabilities of modern LLMS's with respect to concurrent programming.\newline{}
This thesis seeks to investigate modern capabilities of LLM's with respect to concurrency. The project intends to develop a benchmarking and evaluation tool for investigating LLM performance on concurrency problems. The benchmark will evaluate both conventional LLM models and agentic coding harnesses, in an attempt to better reflect modern LLM usage in software development. A potential extension would be to extend the benchmark into a development tool that can provide feedback to AI tools by iteratively evaluating their implmentation and guiding them towards safer concurrent code.
% \section*{From the thesis topic document}
% LLMs are a useful tool for writing code. However, programs generated by LLMs are not guaranteed to be bug-free. Concurrency bugs are specially difficult to identify, as they require reasoning about multiple execution interleavings. Luckily, there exist tools that can automatically find certain types of concurrency errors. In this project, the goal is to use LLMs to generate thread-safe data structures and evaluate whether they contain concurrency bugs. \cite{huang2026concur}

\printbibliography
\end{document}
